% =====================================================================
% contribuciones.tex
% Responsabilidad única: declarar las fuentes de datos, la contribución de
% cada integrante y el uso de herramientas de inteligencia artificial, como
% pide el apartado 8 del enunciado. Sin cifras.
% =====================================================================
\section{Datos, contribuciones y uso de inteligencia artificial}
\label{sec:contribuciones}

\paragraph{Datos.} Las etapas 01 a 07 usan una base simulada que genera
\codigo{proyecto/datos/simulacion.py} con la semilla del proyecto, así que
cualquiera puede reproducirlas. La etapa 08 usa la exportación del sistema
del laboratorio entre junio de 2024 y agosto de 2026, con los clientes, los
puntos de muestreo y los analistas reemplazados por códigos aleatorios. Esa
exportación es confidencial y no viaja con la entrega. En su lugar la
carpeta \codigo{datos/} trae cinco filas con todos los códigos y textos
libres reemplazados y el diccionario de las 36 columnas (tipo, descripción,
papel en el modelo y momento en que se conoce cada una). Sin el archivo la
etapa 08 se ejecuta igual y muestra una nota en lugar de cifras.

\paragraph{Contribuciones.} Los tres integrantes trabajamos de forma
equitativa en todo el proyecto: La formulación de las preguntas, el diseño
del protocolo, el código, la lectura de los resultados y la redacción. Daniel
Arias Mutis, que trabaja en el laboratorio, consiguió y anonimizó la
exportación real y propuso la idea de los lotes de Poisson y de los segmentos
por volumen.

\paragraph{Uso de inteligencia artificial.} Usamos Claude (Anthropic) como
asistente de programación y de redacción: Para escribir y reorganizar
módulos del paquete, armar los guiones de las etapas y revisar el texto. No
le dejamos ninguna cifra a la memoria del asistente ni a la nuestra. Todas
las cifras del informe salen de los artefactos que escriben los notebooks
(macros de \codigo{informe/generado/}) y \codigo{verificar.py} comprueba,
entre otras cosas, la regla de capas, que no haya fugas por momento, que solo
las etapas 07 y 08 lean la prueba y que el texto no tenga cifras escritas a
mano. Además revisamos a mano las decisiones de cada etapa y los resultados
que no nos cuadraban antes de darlos por buenos.
